\subsection{Environment and Agent Interface}\label{sec:environment}
\subsubsection{State and transition model}
Let $\mathcal{G}=\{0,\ldots,W-1\}\times\{0,\ldots,H-1\}$ be the grid and
$\mathcal{I}_t$ the living animals at step $t$. World state $s_t$ includes their
positions, headings, energy, health, age, sex labels, reproductive clocks and
ancestry; resource quantities; obstacles; terrain; cumulative counters; and
random-generator state. Each animal $i$ has species $z_i$, private controller
parameters $\theta_i$, and training state $\xi_i$. The simulator is naturally a
partially observed stochastic game with a variable agent set:
\begin{equation}
 o_{i,t}=O_i(s_t),\qquad a_{i,t}\sim\pi_{\theta_i}(\cdot\mid o_{i,t}),\qquad
 (s_{t+1},r_t,\mathcal{I}_{t+1})=T(s_t,a_t;\omega_t),
\end{equation}
where $\omega_t$ denotes draws from the world's explicit RNG. An experiment
consists of finite episodes, called \emph{runs} by the CLI. Independent learners
and changing opponents make each learner's experienced transition distribution
nonstationary.

Species names are arbitrary configuration keys. Diet determines feeding role;
exact species identity determines mating and species-specific reporting.
Several species can share a brain template without sharing live parameters.
The current reproduction model uses two reproductive sex labels; it does not
model a complete biological account of sex or development.

\subsubsection{Physical update order}
All actions are collected before the physical transition. The world resolves
turning and voluntary movement, applies river drift, resolves feeding and
simultaneous attacks, charges living/action demand, applies starvation or
water-related harm, records deaths, resolves eligible births, regrows resources
and constructs observations. Parent controllers receive the completed step's
learning update before newborn controllers are mixed from them. This ordering
matters: an animal can complete its feeding action in the step it is killed,
and a newborn has no experience transition on its birth step.

Initial occupancy blocks moves; swaps and following into a cell vacated in the
same step are disallowed. Seeded contention resolves competing empty-cell
claims. Resource eaters feed at their current or forward cell, bounded by
available resources and remaining energy capacity. Prey eaters attack a
resource-eating target in their forward cell. Damage credit is capped at actual
remaining target health; exactly one seeded attacker receives kill energy.

Energy and health have distinct roles. Living and action demand spend energy;
exhausted animals receive starvation damage. Ordinary starvation/combat health
has no recovery mechanic, and dead bodies do not persist as edible corpses.
Resource regrowth is capped. Newborn initial energy is an explicit grant, so the
system should not be treated as a closed energy-conserving ecosystem.

\subsubsection{Actions, terrain and visibility}
Species configure their discrete action list. The available vocabulary includes
forward motion, turns, rest, feeding, attack, reproduction and, on appropriate
maps, swimming. An environment need not include irrelevant outputs: dry presets
can omit swimming. Policies may choose infeasible actions and receive their
physical cost/failure feedback; maskable PPO restricts some immediately
infeasible choices.

Terrain adds water depth, current vectors and elevation. Entering water requires
a swim action and the configured health/energy eligibility; swimming and uphill
motion can incur additional demand. Currents accumulate fractional displacement
and traverse intermediate cells under terrain constraints. There is no
separate hand-coded route-planning or river-crossing policy. Geography presets
include a straight river, highlands, enclosed valleys, a meandering floodplain,
a branching delta, alpine passes and lake islands.

Sight uses a circular base radius and integer line-of-sight tests with wall
occlusion. Heading, other animals and elevation do not occlude sight. A map with
water can extend range directionally. Writing $R_i$ for base radius, $B_i$ for
the maximum water bonus and $w(p_i,p)$ for counted intermediate water cells,
the range condition is
\begin{equation}
 \lVert p-p_i\rVert_2\leq R_i+\min\{B_i,w(p_i,p)\},
\end{equation}
combined with the ordinary wall-visibility check. This is an authored sensing
rule, not an empirical model of animal vision.

\subsubsection{Observation structure}
Observations contain eight self/local-movement values and up to $k$ visible
targets in each of three diet-based categories: resource patches,
resource-eating animals and prey-eating animals. Each target slot holds
$(\Delta x,\Delta y,d,\text{presence})$. Coordinates/distances are normalized by
the maximum query radius. Targets are ordered by distance, with authoritative
insertion-order ties, after visibility filtering; missing slots are zero padded.
No target identity, health, sex or brain state is disclosed by those slots.

Versions 2 and 3 add 12 local terrain values; version 3 adds ten mating-related
values, including personal readiness and one nearest-visible-mate summary.
Thus the input dimension is
\begin{equation}\label{eq:obswidth}
 d(v,k)=8+12k+12\ind[v\geq2]+10\ind[v=3],\qquad 1\leq k\leq16.
\end{equation}
For $k=1$ the widths are 20, 32 and 42; for $k=4$ they are 56, 68 and 78.
Terrain fields are padded on flat maps where needed. Mate information remains a
single summary rather than $k$ extra mates. Readiness does not expose future
consent, and observing never consumes reproductive RNG draws.

\subsubsection{Feedback and transition boundaries}
For animal $i$, write $F_i$ for actual energy gained, $D_i^+$ and $D_i^-$ for
actual damage dealt/taken, $C_i$ for energy demand, $M_i$ for a failed attempted
action, $B_i$ for successful-parent credit and $Z_i$ for death. The environmental
reward is described by
\begin{equation}\label{eq:reward}
 r_i=c_FF_i+c_+D_i^++c_-D_i^-+c_CC_i+c_MM_i+c_BB_i+c_ZZ_i.
\end{equation}
Coefficients belong to species JSON, not map-name branches. This expression
summarizes accounting; it is not a change to numerical summation order. Demand
is counted before energy is floored at zero. Physical spending and scalar
feedback do not debit energy twice. Failed reproduction does not pay the
successful-birth cost or start a cooldown, although normal living demand and
failure feedback still apply. Survival and proximity have no implicit bonus.

Death terminates an individual's transition. The episode time limit truncates
it; value-based and actor--critic methods retain their bootstrap semantics.
Actor-only finite-return methods keep their own return definitions. Recorded
reward remains individual; MAPPO intentionally uses a species-mean adaptation
for training. Environmental feedback is not itself the learner's loss.
